\documentclass[a4paper,12pt]{article}

% ============================================================
% Общий преамбул для отчётов по лабораторным работам (labs_lpr)
% Подключается из report/main.tex командой:
%   % ============================================================
% Общий преамбул для отчётов по лабораторным работам (labs_lpr)
% Подключается из report/main.tex командой:
%   % ============================================================
% Общий преамбул для отчётов по лабораторным работам (labs_lpr)
% Подключается из report/main.tex командой:
%   \input{../../common/latex_template/preamble.tex}
% ============================================================

% Поддержка русского языка и кодировки
\usepackage[T2A]{fontenc}
\usepackage{lmodern}
\usepackage[utf8]{inputenc}
\usepackage[russian]{babel}

% Математические пакеты
\usepackage{amsmath, amssymb}

% Настройка полей
\usepackage{geometry}
\geometry{top=2cm, bottom=2cm, left=2.5cm, right=1.5cm}

% Колонтитулы
\usepackage{fancyhdr}
\usepackage{graphicx}

% Современные таблицы
\usepackage{tabularray}

% Работа с csv файлами
\usepackage{csvsimple-l3}
\usepackage{tikz}

% Графики
\usepackage{pgfplots}
\pgfplotsset{compat=1.18}

% Ссылки
% \usepackage{hyperref} % гипперссылки, нужно оформление вместо красного квадрата
\usepackage[russian]{cleveref} % умные ссылки, нужно переоформление для формул, рисунков и тд (вместо ф-л что-то нужно придумать)

% ============================================================

% Поддержка русского языка и кодировки
\usepackage[T2A]{fontenc}
\usepackage{lmodern}
\usepackage[utf8]{inputenc}
\usepackage[russian]{babel}

% Математические пакеты
\usepackage{amsmath, amssymb}

% Настройка полей
\usepackage{geometry}
\geometry{top=2cm, bottom=2cm, left=2.5cm, right=1.5cm}

% Колонтитулы
\usepackage{fancyhdr}
\usepackage{graphicx}

% Современные таблицы
\usepackage{tabularray}

% Работа с csv файлами
\usepackage{csvsimple-l3}
\usepackage{tikz}

% Графики
\usepackage{pgfplots}
\pgfplotsset{compat=1.18}

% Ссылки
% \usepackage{hyperref} % гипперссылки, нужно оформление вместо красного квадрата
\usepackage[russian]{cleveref} % умные ссылки, нужно переоформление для формул, рисунков и тд (вместо ф-л что-то нужно придумать)

% ============================================================

% Поддержка русского языка и кодировки
\usepackage[T2A]{fontenc}
\usepackage{lmodern}
\usepackage[utf8]{inputenc}
\usepackage[russian]{babel}

% Математические пакеты
\usepackage{amsmath, amssymb}

% Настройка полей
\usepackage{geometry}
\geometry{top=2cm, bottom=2cm, left=2.5cm, right=1.5cm}

% Колонтитулы
\usepackage{fancyhdr}
\usepackage{graphicx}

% Современные таблицы
\usepackage{tabularray}

% Работа с csv файлами
\usepackage{csvsimple-l3}
\usepackage{tikz}

% Графики
\usepackage{pgfplots}
\pgfplotsset{compat=1.18}

% Ссылки
% \usepackage{hyperref} % гипперссылки, нужно оформление вместо красного квадрата
\usepackage[russian]{cleveref} % умные ссылки, нужно переоформление для формул, рисунков и тд (вместо ф-л что-то нужно придумать)


\usepackage{subcaption}
\usepackage{circuitikz}
\usepackage{float}

% Числа на осях: десятичная запятая, без разделителя тысяч
\pgfplotsset{/pgf/number format/.cd, use comma, 1000 sep={}}

% --- Служебные макросы ------------------------------------------------------
\newcommand{\qq}{\textbf{???}}
\newcommand{\fillme}{\textbf{?? НЕОБХОДИМО ЗАПОЛНИТЬ ??}}
\newcommand{\filename}[1]{\texttt{\detokenize{#1}}}
% Итоговые числа из data/processed/results.tex; нет значения -> ???
\IfFileExists{../data/processed/results.tex}{% Автоматически создано scripts/process_data.py -- вручную не править.
\expandafter\def\csname res@Czero\endcsname{1{,}14}
\expandafter\def\csname res@criticalTheory\endcsname{7480}
\expandafter\def\csname res@criticalObserved\endcsname{6520}
\expandafter\def\csname res@free1\endcsname{8{,}53}
\expandafter\def\csname res@free2\endcsname{1{,}92}
\expandafter\def\csname res@calc1\endcsname{8{,}64}
\expandafter\def\csname res@peakFreq1\endcsname{6440}
\expandafter\def\csname res@peakV1\endcsname{20{,}0}
\expandafter\def\csname res@ach1\endcsname{8{,}51}
\expandafter\def\csname res@phase1\endcsname{8{,}45}
\expandafter\def\csname res@calc2\endcsname{1{,}85}
\expandafter\def\csname res@hasperiod\endcsname{1}
\expandafter\def\csname res@hascritical\endcsname{1}
\expandafter\def\csname res@hasamp1\endcsname{1}
\expandafter\def\csname res@hasph1\endcsname{1}
\expandafter\def\csname res@periodSlope\endcsname{0{,}995 \pm 0{,}004}
\expandafter\def\csname res@periodIntercept\endcsname{0{,}4 \pm 0{,}6}
\expandafter\def\csname res@criticalSlope\endcsname{(1{,}34 \pm 0{,}09) \cdot 10^{6}}
\expandafter\def\csname res@criticalIntercept\endcsname{0{,}25 \pm 0{,}22}
\expandafter\def\csname res@criticalFit\endcsname{7260 \pm 250}
\expandafter\def\csname res@hasenv2496\endcsname{1}
\expandafter\def\csname res@hasenv2499\endcsname{1}
\expandafter\def\csname res@hasenv2500\endcsname{1}
\expandafter\def\csname res@hasenv2502\endcsname{1}
\expandafter\def\csname res@hasenv2503\endcsname{1}
\expandafter\def\csname res@rise1\endcsname{8{,}99 \pm 0{,}11}
\expandafter\def\csname res@photoTheta\endcsname{0{,}3567 \pm 0{,}0027}
\expandafter\def\csname res@fall1\endcsname{8{,}81 \pm 0{,}07}
}{}
\newcommand{\res}[1]{\ifcsname res@#1\endcsname\ensuremath{\csname res@#1\endcsname}\else\qq\fi}
% Фото/зарисовка из report/img/. Если файла ещё нет -- рамка-заглушка.
\newcommand{\labimage}[2]{%
    \IfFileExists{img/#1}{%
        \includegraphics[width=#2\linewidth]{img/#1}%
    }{%
        \fbox{\parbox[c][0.6\dimexpr#2\linewidth\relax][c]{%
                \dimexpr#2\linewidth-2\fboxsep-2\fboxrule\relax}{%
            \centering\small\filename{img/#1}\\[2mm]Материал будет добавлен}}%
    }%
}
\newcommand{\vx}{\text{вх}}
\newcommand{\vy}{\text{вых}}

% Колонтитулы (как в остальных работах репозитория)
\pagestyle{fancy}
\fancyhf{}
\fancyhead[L]{%
    \IfFileExists{../../common/latex_template/miptlogo.png}{%
        \includegraphics[height=1cm]{../../common/latex_template/miptlogo.png}%
    }{%
        \textbf{[ЛОГОТИП МФТИ]}%
    }%
}
\fancyhead[R]{\textbf{Емельянов Дмитрий}}
\fancyfoot[C]{\thepage}
\renewcommand{\headrulewidth}{0.4pt}
\setlength{\headheight}{33pt}

\begin{document}

% --- ТИТУЛЬНЫЙ ЛИСТ ---
\begin{titlepage}
\begin{center}
    {\large МОСКОВСКИЙ ФИЗИКО-ТЕХНИЧЕСКИЙ ИНСТИТУТ}
\end{center}
\begin{center}
    {\large Физтех-школа физики и исследований им. Ландау\break ЛФИ}
\end{center}

    \vspace{3cm}

\vspace{0.15cm}
{
\begin{center}
    \huge \bf Лабораторная работа №3.2.5 по курсу Общая физика\\
    \vspace{0.4cm}
    \LARGE Свободные и вынужденные колебания в электрическом контуре
\end{center}
}
\vspace{1.5cm}
\begin{flushright}
Емельянов Дмитрий\\
\vspace{0.4cm}
Б02-502
\end{flushright}
\vfill
\begin{center}
    \large 09.11.2026, г. Долгопрудный
\end{center}
\end{titlepage}

\newcommand{\panel}[2]{\begin{minipage}[t]{.47\linewidth}\centering\labimage{#1}{.92}\\[2pt]\small #2\end{minipage}}
\section{Аннотация}
Цель работы --- исследовать свободный отклик электрического контура и его отклик на синусоидальное возбуждение. По периоду при разных ёмкостях проверяется формула Томсона. По изменению затухания с сопротивлением определяется граница колебательного режима. Для двух сопротивлений сравниваются оценки добротности по параметрам цепи, свободному затуханию, резонансным кривым и переходным процессам.
\section{Теоретические сведения}
Пусть $C=C_{\textrm{маг}}+C_0$ --- полная ёмкость, $R_\Sigma=R+R_L$ --- сумма сопротивлений магазина и обмотки. Между импульсами конденсатор разряжается через катушку и резистор. Из закона Кирхгофа для заряда $q=CU_C$ следует
\[
 \ddot q+2\gamma\dot q+\omega_0^2q=0,\qquad
 \gamma=R_\Sigma/(2L),\quad \omega_0=(LC)^{-1/2}.
\]
При $\gamma<\omega_0$ напряжение равно $Ae^{-\gamma t}\cos(\omega_dt+\alpha)$, где $\omega_d=\sqrt{\omega_0^2-\gamma^2}$. Период $T_d=2\pi/\omega_d$ при слабом затухании близок к $T_0=2\pi\sqrt{LC}$. При $R_\Sigma=R_{\textrm{кр}}=2\sqrt{L/C}$ корни характеристического уравнения совпадают, при большем сопротивлении разряд апериодический.
\[
 \Theta=\gamma T_d=\frac1n\ln\frac{|U_k|}{|U_{k+n}|},
 \qquad Q=\frac{\omega_0}{2\gamma}=\frac{1}{R_\Sigma}\sqrt{\frac LC}.
\]
Величина $\pi/\Theta=\omega_d/(2\gamma)$ приближает $Q$ при слабом затухании. На фазовой плоскости $X=U_C$, $Y=RI\propto\dot U_C$ траектория образует спираль. Одноимённые пересечения одного луча отстоят на целое число периодов, поэтому отношение их расстояний от центра даёт тот же декремент.
\par При гармоническом возбуждении в модели осциллятора с постоянной амплитудой силы правая часть уравнения пропорциональна $\cos\omega t$. Амплитуда и фаза установившегося отклика имеют вид
\[
 U_C\propto[(\omega_0^2-\omega^2)^2+4\gamma^2\omega^2]^{-1/2},
 \quad \varphi=-\operatorname{atan2}(2\gamma\omega,\omega_0^2-\omega^2).
\]
Максимум амплитуды приходится на $\sqrt{\omega_0^2-2\gamma^2}$, а запаздывание на четверть периода --- на $\omega_0$. При включении цуга на резонансе в приближении слабого затухания огибающая растёт как $U_0(1-e^{-\gamma t})$, после выключения спадает как $U_{\textrm{кон}}e^{-\gamma t}$.
\section{Оборудование и подключение}
Генератор АКИП-3409/4 через соединительную коробку с ёмкостью связи $C_1$ возбуждает контур. Канал 1 осциллографа АКТАКОМ ADS-6142H измеряет $U_C$. Для фазового портрета канал 2 измеряет напряжение на магазине МСР-60, для резонансных кривых --- выход генератора через коаксиальный разветвитель. Ёмкость задаётся магазином Р5025, индуктивность --- Р567 типа МИСП. Активное сопротивление катушки и её индуктивность измеряются LCR-метром.
\par Инструментальные погрешности курсоров времени и напряжения, частоты генератора, магазинов и LCR-метра: \qq. Для каждого диапазона следует записать цену деления, разрешение и паспортную точность.
\begin{figure}[H]\centering\panel{setup.jpg}{Общий вид установки}\hfill\panel{free.png}{Свободный разряд после импульса}\caption{Подключение приборов и сигнал канала 1.}\end{figure}

\section{Методика измерений}
\begin{enumerate}
\item[2.1.] Подать на канал 1 импульсы длительностью 10~мкс, амплитудой 20~В с частотой 100~Гц. Настроить запуск по каналу 1, затем подключить контур.
\item[2.2.] При $L=100$~мГн и $R=0$ записать 5--10 значений периода для ёмкостей магазина от 0 до 9~нФ. Для каждой точки измерить интервал $\Delta t$ между одноимёнными максимумами через $n$ периодов. Нулевую установку измерить отдельно для определения $C_0$.
\item[2.3.] Для частоты 6,5~кГц расчётная полная ёмкость $C^*=[L(2\pi\nu_0)^2]^{-1}\approx6{,}00$~нФ, $R_{\textrm{кр}}\approx8{,}17$~кОм. Подобрать установку магазина с учётом $C_0$, записать фактическое значение. Повышая $R$, зафиксировать исчезновение колебаний. Для 6--8 сопротивлений в диапазоне $(0{,}05$--$0{,}25)R_{\textrm{кр}}$ измерить амплитуды двух одноимённых пиков и число периодов между ними. Сохранить крайние настройки $R_1,R_2$ для последующих серий.
\item[2.4.] Подать $U_R$ на канал 2 и наблюдать оба сигнала во времени. При частоте повторения около 400--500~Гц перейти в режим XY. При неизменной ёмкости проследить изменение спирали от $R_1$ до $R_2$. Для обеих настроек снять координаты пересечений одного луча через $n$ оборотов.
\item[2.5.] Перейти к синусоидальному возбуждению и временной развёртке. Совместить нулевые уровни каналов. При постоянной амплитуде генератора найти максимум $U_C$ для каждого $R$. Снять минимум по десять точек с каждой стороны максимума вплоть до $0{,}4U_{\textrm{рез}}$. Записывать частоту, размах $2U_C$ и запаздывание $\Delta t$ между одноимёнными переходами через ноль. Фаза выхода $\varphi=-2\pi\nu\Delta t$.
\item[2.6.] На резонансе подать цуги по 15 циклов с периодом повторения 20~мс. Для каждого сопротивления взять по 3--4 пары пиков на участках нарастания и затухания, а также установившуюся амплитуду $U_0$. Если плато не достигнуто, измерить $U_0$ при непрерывном возбуждении той же амплитуды. При $R_1$ отстроить частоту и записать биения. После отключения установки измерить $R_L$ и $L$ на частотах 50, 500 и 1500~Гц.
\end{enumerate}
\begin{figure}[H]\centering\panel{spiral.png}{Фазовая плоскость $(U_C,U_R)$}\hfill\panel{phase.png}{Сигнал генератора и отклик контура}\caption{Осциллограммы для измерения декремента и сдвига фаз.}\end{figure}
\section{Результаты и обработка}
\begin{table}[H]\centering\small\caption{Периоды свободного разряда при $R=0$ по лабораторному журналу.}
\begin{tabular}{l|rrrrrrr}
$C_{\textrm{маг}}$, нФ &0&1&2&3&5&7&9\\
$T$, мкс &67,2&91,8&112,0&127,0&155,1&178,8&199,7
\end{tabular}\end{table}
\begin{table}[H]\centering\small\caption{Последовательные амплитуды свободного разряда. Единица отсчёта в журнале не указана.}
\begin{tabular}{l|rrrrrrr}
$R$, Ом&410&600&800&1000&1200&1600&2000\\
$U_1$&580&660&470&610&360&550&524\\
$U_2$&400&371&250&248&150&140&104\\
$U_3$&284&219&140&110&70&41&20\\
$U_4$&192&130&80&38&---&---&---
\end{tabular}\end{table}
\paragraph{Периоды (п. 2.7.1).} $T_{\textrm{эксп}}=\Delta t/n$, $C_0=T_{\textrm{эксп}}^2(0)/(4\pi^2 L)=\res{Czero}$~нФ. По семи точкам методом наименьших квадратов получена прямая $T_{\textrm{эксп}}=kT_0+b$, где $k=\res{periodSlope}$ и $b=\res{periodIntercept}$~мкс. Для формулы Томсона ожидаются $k=1$ и $b=0$. Для независимых отсчётов $u(T_{\textrm{эксп}})=u(\Delta t)/n$. При подстановке $C_0$ удобно считать $T_0=\sqrt{T^2(0)+4\pi^2 LC_{\textrm{маг}}}$, чтобы учесть общий параметр $L$ один раз. Нулевая точка задаёт калибровку.
\paragraph{Критическое сопротивление (пп. 2.7.2--3).} Для свободного разряда построить $Y=1/\Theta^2$ от $X=1/R_\Sigma^2$. Из уравнения контура $Y=aX-1/(4\pi^2)$, $a=R_{\textrm{кр}}^2/(4\pi^2)$. По наклону линейного участка у начала координат $R_{\textrm{кр}}=2\pi\sqrt a$, $u(R_{\textrm{кр}})=\pi u(a)/\sqrt a$. Свободная МНК-аппроксимация $Y=aX+b$ по семи точкам даёт $a=\res{criticalSlope}$~Ом$^2$ и $b=\res{criticalIntercept}$. Погрешности параметров обеих прямых оценены по разбросу остатков. Результат по наклону: \res{criticalFit}~Ом; по параметрам $L,C$: \res{criticalTheory}~Ом; по границе режимов $R_{\textrm{переход}}+R_L$: \res{criticalObserved}~Ом.
\paragraph{Свободный разряд и спираль (пп. 2.7.4--6).} По крайним декрементам и координатам витков вычислить $Q=\pi/\Theta$. Для независимых амплитуд\[u(\Theta)=\sqrt{[u(U_k)/U_k]^2+[u(U_{k+n})/U_{k+n}]^2}/n.\]Для сравнения вычислить $Q$ по $L,C,R_\Sigma$ с указанием частоты измерения катушки.
\paragraph{Резонансные кривые (пп. 2.7.7--10).} Совместить АЧХ для двух сопротивлений в координатах $U_C/U_{\textrm{рез}}$ и $\nu/\nu_{\textrm{рез}}$. Линейно интерполировать частоты пересечения уровня $1/\sqrt2$ на обоих склонах. Оценка при узком резонансе $Q_{\textrm{А}}=\nu_{\textrm{рез}}/(\nu_+-\nu_-)$. Построить обе ФЧХ в координатах $(\omega,\varphi)$. Отразить часть ниже $-\pi/2$ по правилу $\varphi'=-\pi-\varphi$. Ширина между пересечениями уровня $-\pi/4$ равна разности частот исходной кривой при $-3\pi/4$ и $-\pi/4$. Тогда $Q_{\varphi}=\omega_{\varphi=-\pi/2}/\Delta\omega$.
\paragraph{Оцифровка видео.} Для $R_1=410$~Ом использована запись IMG\_2494.MOV. После выравнивания перспективы обоих дисплеев настройки генератора разделены по изменению выделенного разряда (10~Гц). Опорные показания прочитаны на дисплее, промежуточные распознаны с учётом монотонного вращения ручки и периода синей трассы. Для 168 различимых настроек сохранено по одному кадру. Из-за размытия младшего разряда частота отдельных точек может отличаться на один-два шага.
\par Для каждого цвета по пикселям трассы построена синусоида. Размах жёлтой синусоиды делится на наибольший размах видеосерии ($\res{peakV1}$~В при $\res{peakFreq1}$~Гц). Время отсчитывается вправо: $\Delta t=t_{\textrm{жёлт}}-t_{\textrm{син}}$ для ближайших максимумов, $\varphi=-2\pi\nu\Delta t$. Жёлтые максимумы на записи расположены левее синих, поэтому измеренная фаза положительна. На этой ветви берём ширину между уровнями $3\pi/4$ и $\pi/4$, центральную частоту --- на уровне $\pi/2$. Это те же уровни относительно середины фазового перехода, что и в методике. Для второго сопротивления в записи отдельного частотного прохода нет.
\paragraph{Цуги и биения (п. 2.7.11).} На участке роста использовать
$\Theta=n^{-1}\ln[(U_0-U_k)/(U_0-U_{k+n})]$, на спаде --- формулу свободного разряда. Сопоставить оценки $\pi/\Theta$. При отстройке свободная составляющая на частоте $\nu_d$ складывается с вынужденной на частоте $\nu$. Ожидаемый интервал максимумов огибающей $1/|\nu-\nu_d|$ сравнить с видеозаписью, учитывая затухание свободной составляющей.
\begin{figure}[H]\centering\panel{burst.png}{Отклик контура на цуг, $R_1=410$~Ом}\hfill\panel{decay.jpg}{Спад того же сигнала, $R_1=410$~Ом}\caption{Пункт 2.6.5. Слева весь отклик на цуг (IMG\_2496), справа участок спада (IMG\_2500).}\end{figure}
\paragraph{Оценка по фотографии IMG\_2500.} Для каждого цикла жёлтой трассы определён полуразмах по нижнему экстремуму и двум соседним верхним. На участке спада от второго до девятого оцифрованного цикла выполнена линейная аппроксимация $\ln A$ от времени в единицах периода. Масштаб по вертикали сокращается в отношении амплитуд. Получены $\Theta=\res{photoTheta}$ и $Q=\res{fall1}$. Здесь указана только стандартная ошибка наклона по остаткам аппроксимации. Измерение относится к $R_1=410$~Ом. По семи точкам участка роста IMG\_2499 до выхода на плато аппроксимация $A(n)=A_0-Be^{-n\Theta}$ даёт $Q=\res{rise1}$; $n$ измеряется в периодах, $A_0$ определено совместно с $B$ и $\Theta$. Обе оценки погрешности характеризуют разброс относительно модели.
\begin{figure}[H]\centering\panel{growth.jpg}{Нарастание при $R_1=410$~Ом, IMG\_2499}\hfill\panel{burst_fast.jpg}{Отклик на цуг при $R_2=2000$~Ом, IMG\_2502}\\[3mm]
\panel{burst_detail.jpg}{Биения при $R_1=410$~Ом и $\nu\approx5{,}7$~кГц, IMG\_2503}\hfill\panel{generator.jpg}{Настройки генератора, IMG\_2504}
\caption{Осциллограммы и настройки генератора в режиме цугов.}\end{figure}
\paragraph{Погрешности и сопоставление (пп. 2.7.12--13).} Для результата $F$ использовать $u^2(F)=\sum_{ij}(\partial_i F)(\partial_j F)\operatorname{cov}(x_i,x_j)$. При обработке роста брать $U_0$ одним общим аргументом для обеих разностей. Для ширины кривой учитывать ошибки частот и вертикальных отсчётов через локальный наклон, для АЧХ также общий уровень нормировки. Разброс независимых пар учитывать по стандартной ошибке среднего. Относительная погрешность равна $100u(Q)/Q$.
\begin{table}[H]\centering\small\caption{Добротность контура (безразмерная) для двух настроек сопротивления.}
\begin{tabular}{c|ccc|cccc}
 &\multicolumn{3}{c|}{Свободные колебания}&\multicolumn{4}{c}{Вынужденные колебания}\\
$R$ & Расчёт & Разряд & Спираль & АЧХ & ФЧХ & Рост & Спад\\\hline
$R_1$ & \res{calc1} & \res{free1} & \res{spiral1} & \res{ach1} & \res{phase1} & \res{rise1} & \res{fall1} \\
$R_2$ & \res{calc2} & \res{free2} & \res{spiral2} & \res{ach2} & \res{phase2} & \res{rise2} & \res{fall2} \\
\end{tabular}\end{table}
\newcommand{\series}[2]{\ifcsname res@has#1\endcsname
\addplot+[only marks,#2,error bars/.cd,x dir=both,x explicit,y dir=both,y explicit] table[col sep=comma,x=x,y=y,x error=sx,y error=sy]{../data/processed/#1.csv};\fi}
\newcommand{\fitline}[1]{\ifcsname res@#1Slope\endcsname\addplot[thin,no marks,black] table[col sep=comma,x=x,y=y]{../data/processed/#1_fit.csv};\fi}
\newcommand{\chart}[4]{\begin{minipage}{.49\linewidth}\centering\begin{tikzpicture}\begin{axis}[width=\linewidth,height=5cm,xlabel={#1},ylabel={#2},scaled ticks=true]\ifcsname res@has#3\endcsname\series{#3}{blue}#4\else\path[use as bounding box] (axis cs:0,0) rectangle (axis cs:1,1);\node at (rel axis cs:.5,.5){Нет отсчётов};\fi\end{axis}\end{tikzpicture}\end{minipage}}
\begin{figure}[H]\centering
\chart{$T_0$, мкс}{$T_{\textrm{эксп}}$, мкс}{period}{\fitline{period}}
\chart{$1/R_\Sigma^2$, Ом$^{-2}$}{$1/\Theta^2$}{critical}{\fitline{critical}}\\
\chart{$\nu/\nu_{\textrm{рез}}$}{$U/U_{\textrm{рез}}$}{amp1}{\series{amp2}{red}}
\ifcsname res@hasph1\endcsname
\chart{$\omega$, рад/с}{$\varphi$, рад}{ph1}{\series{ph2}{red}}
\else\ifcsname res@hasph2\endcsname
\chart{$\omega$, рад/с}{$\varphi$, рад}{ph2}{}
\else\begin{minipage}{.48\linewidth}\centering\small Фазочастотная характеристика\\[2mm]Для построения нужны пары\\ «частота --- задержка».\end{minipage}\fi\fi
\caption{Периоды, декременты и резонансные кривые. Вверху показаны все семь точек каждой серии и прямые МНК. Внизу показаны АЧХ и ФЧХ по видеозаписи при $R_1=410$~Ом. Фаза отсчитана относительно синего сигнала генератора.}\end{figure}

\newcommand{\envelope}[2]{\ifcsname res@hasenv#1\endcsname
\addplot+[thin,solid,mark size=1.3pt] table[col sep=comma,x=x,y=y]{../data/processed/env#1.csv};\addlegendentry{#2}\fi}
\begin{figure}[H]\centering
\begin{tikzpicture}\begin{axis}[width=.94\linewidth,height=7cm,xlabel={Время от начала отсчёта, периоды},ylabel={$A/A_{\textrm{макс}}$},legend style={font=\small},legend columns=3,legend style={at={(0.5,-0.3)},anchor=north,font=\footnotesize}]
\envelope{2496}{2496, весь цуг}\envelope{2499}{2499, рост}\envelope{2500}{2500, спад}\envelope{2502}{2502, весь цуг}\envelope{2503}{2503, биения}
\end{axis}\end{tikzpicture}
\caption{Огибающие по пяти фотографиям. Линии соединяют отсчёты одной осциллограммы в порядке времени. Каждая точка --- полуразмах одного цикла, нормированный на наибольший полуразмах в том же кадре. Для каждого снимка начало отсчёта соответствует середине первого оцифрованного цикла. Снимок 2502 относится к $R_2$, остальные --- к $R_1$.}\end{figure}

\section{Обсуждение и заключение}
Период при увеличении ёмкости магазина от 0 до 9~нФ вырос от 67,2 до 199,7~мкс. При обработке использованы параметры катушки на частоте 1500~Гц ($L=100{,}04$~мГн, $R_L=23{,}110$~Ом) --- ближайшей к рабочей из трёх измеренных. Для свободного затухания ёмкость магазина равнялась 6~нФ, полная ёмкость около 7,14~нФ. При увеличении $R$ от 410 до 2000~Ом добротность по декременту уменьшилась с $\res{free1}$ до $\res{free2}$. Значение критического сопротивления по прямой, построенной по семи точкам, ближе к расчётному значению, чем визуальная оценка перехода.
\par Перед снятием резонансных кривых ёмкость магазина уменьшили до 5~нФ. При $R_1=410$~Ом зарегистрирован максимум на 6420~Гц с размахом 20,2~В. Журнал указывает диапазон видеосъёмки 5500--7500~Гц. По видео максимум размаха составляет $\res{peakV1}$~В при $\res{peakFreq1}$~Гц. Ширины амплитудной и фазовой характеристик дают близкие оценки добротности $Q_{\textrm{АЧХ}}=\res{ach1}$ и $Q_{\textrm{ФЧХ}}=\res{phase1}$. Для величин из журнала инструментальные погрешности пока не установлены.
\end{document}
